\section{Nettverkslaget -- Oversikt (4.1)}
%=============================================================================

\subsection{Nettverkslagets rolle}

\begin{keybox}[Nettverkslaget -- oversikt]
Nettverkslaget transporterer \textbf{segmenter} fra sendende vert til mottakende vert.
\begin{itemize}[noitemsep]
    \item Sendersiden: kapsler inn segmenter i \textbf{datagrammer} (pakker)
    \item Mottakersiden: leverer segmenter til transportlaget
    \item Nettverkslagsprotokoller finnes i \emph{alle} verter og rutere
    \item Ruteren unders\o ker header-felt i alle IP-datagrammer som passerer gjennom den
\end{itemize}
\end{keybox}

\subsection{Forwarding vs. Routing}

\begin{exambox}[Eksamensrelevant: Forwarding vs. Routing (V2025, V2023)]
\textbf{Forwarding} og \textbf{routing} er to sentrale begreper som ofte forveksles p\aa\ eksamen:
\begin{itemize}[noitemsep]
    \item \textbf{Forwarding (videresending):} Flytte pakker fra ruterens inngangsport til riktig utgangsport. Dette er en \emph{lokal} handling i \emph{hver enkelt ruter} (dataplanet).
    \item \textbf{Routing (ruting):} Bestemme ruten pakker skal ta fra kilde til destinasjon. Dette er en \emph{global}, \emph{nettverksomfattende} prosess (kontrollplanet).
\end{itemize}

\medskip
\textbf{Analogi:} Forwarding = ta riktig avkj\o ring i et veikryss. Routing = planlegge hele reiseruten fra Oslo til Trondheim.

\medskip
\textbf{Fra eksamen V2025:} ``What is the difference between forwarding and routing?''
\begin{itemize}[noitemsep]
    \item Forwarding: move packets from router's input to appropriate router output (local action, data plane)
    \item Routing: determine route taken by packets from source to destination (global action, control plane)
\end{itemize}
\end{exambox}

\subsection{Dataplanet vs. Kontrollplanet}

\begin{keybox}[To plan i nettverkslaget]
\begin{itemize}[noitemsep]
    \item \textbf{Dataplanet:}
    \begin{itemize}[noitemsep]
        \item Lokal, per-ruter funksjon
        \item Bestemmer hvordan datagrammer som ankommer inngangsport videresendes til utgangsport
        \item Forwarding-funksjon -- opererer i \textbf{nanosekunders} tidsramme (hardware)
    \end{itemize}
    \item \textbf{Kontrollplanet:}
    \begin{itemize}[noitemsep]
        \item Nettverksomfattende logikk
        \item Bestemmer hvordan datagrammer rutes mellom rutere langs ende-til-ende-stier
        \item Opererer i \textbf{millisekunders} tidsramme (software)
        \item To tilnerminger: tradisjonelle rutingsalgoritmer (per-ruter) eller SDN (Software-Defined Networking)
    \end{itemize}
\end{itemize}
\end{keybox}

\subsection{Per-ruter kontrollplan vs. SDN}

\begin{keybox}[To kontrollplan-tilnerminger]
\textbf{Tradisjonelt (per-ruter kontrollplan):}
\begin{itemize}[noitemsep]
    \item Hver ruter har sin egen rutingsalgoritme
    \item Rutere kommuniserer seg imellom for \aa\ bygge forwarding-tabeller
    \item Rutingsalgoritmen kj\o rer i \emph{hver} ruter
\end{itemize}

\textbf{SDN (Software-Defined Networking):}
\begin{itemize}[noitemsep]
    \item En ekstern (fjern) kontroller beregner og distribuerer forwarding-tabeller
    \item Rutere/svitsjer f\aa r sine flow-tabeller fra den sentrale kontrolleren
    \item Separerer kontrollplanet fra dataplanet fysisk
\end{itemize}
\end{keybox}

\subsection{Nettverkstjenestemodellen}

\begin{keybox}[Best Effort-tjeneste]
Internettets nettverkslag tilbyr kun \textbf{best effort}-tjeneste:
\begin{itemize}[noitemsep]
    \item \textbf{Ingen} garanti for b\aa ndbredde
    \item \textbf{Ingen} garanti mot pakketap
    \item \textbf{Ingen} garanti for rekkef\o lge
    \item \textbf{Ingen} garanti for forsinkelse (timing)
    \item \textbf{Ingen} garanti for minimal forsinkelse mellom pakker (jitter)
\end{itemize}
\end{keybox}

\begin{warnbox}[Sammenligning med andre nettverksarkitekturer]
Andre nettverksarkitekturer (som ATM) tilb\o d sterkere garantier:
\begin{itemize}[noitemsep]
    \item \textbf{ATM CBR} (Constant Bit Rate): garantert b\aa ndbredde, rekkef\o lge, timing, ingen tap
    \item \textbf{ATM ABR} (Available Bit Rate): garantert minimumsb\aa ndbredde, ingen tap
\end{itemize}
Internettets ``best effort'' har likevel vunnet fordi enkelheten tillater innovasjon i ytterkantene.
\end{warnbox}

%=============================================================================
\section{Hva er inne i en ruter? (4.2)}
%=============================================================================

\subsection{Ruterarkitektur -- oversikt}

\begin{keybox}[Generisk ruterarkitektur]
En ruter best\aa r av fire hovedkomponenter:
\begin{enumerate}[noitemsep]
    \item \textbf{Inngangsporter (Input ports):} Fysisk mottak, linklags-protokoll, oppslag i forwarding-tabell, k\o
    \item \textbf{Svitsjefabrikk (Switching fabric):} Kobler inngangsporter til utgangsporter
    \item \textbf{Utgangsporter (Output ports):} Buffring, k\o disiplin, linklags-protokoll, fysisk sending
    \item \textbf{Rutingprosessor:} Kj\o rer rutingsprotokoller, administrasjon (kontrollplanet)
\end{enumerate}

\medskip
\textbf{Viktig skille:}
\begin{itemize}[noitemsep]
    \item Kontrollplanet (rutingprosessor): software, millisekunder
    \item Dataplanet (inngangsporter, svitsjefabrikk, utgangsporter): hardware, nanosekunder
\end{itemize}
\end{keybox}

\subsection{Inngangsport-funksjoner}

\begin{keybox}[Inngangsportens tre steg]
\begin{enumerate}[noitemsep]
    \item \textbf{Linjeterminering:} Fysisk lag -- bit-niv\aa\ mottak
    \item \textbf{Linklags-protokoll:} F.eks. Ethernet (motta ramme)
    \item \textbf{Oppslag og forwarding:} Bruk header-verdier til \aa\ sl\aa\ opp utgangsport i forwarding-tabellen (``match plus action'')
\end{enumerate}

\medskip
\textbf{Desentralisert svitsjing:}
\begin{itemize}[noitemsep]
    \item Forwarding-tabellen er kopiert til inngangsportens minne
    \item M\aa l: fullf\o re input-port-prosessering med \textbf{linjehastighet} (line speed)
    \item \textbf{Input port queuing:} Oppst\aa r n\aa r datagrammer ankommer raskere enn switching-hastigheten
\end{itemize}
\end{keybox}

\subsubsection{Destinasjonsbasert vs. generalisert forwarding}
\begin{itemize}[noitemsep]
    \item \textbf{Destinasjonsbasert forwarding:} Videresending basert kun p\aa\ destinasjons-IP-adresse (tradisjonelt)
    \item \textbf{Generalisert forwarding:} Videresending basert p\aa\ et vilk\aa rlig sett av header-verdier (SDN/OpenFlow)
\end{itemize}

\subsection{Longest Prefix Match}

\begin{exambox}[Eksamensrelevant: Longest Prefix Match (V2023 Q4)]
\textbf{Longest prefix match:} N\aa r man sl\aa r opp i forwarding-tabellen for en gitt destinasjonsadresse, bruk den oppf\o ringen med \emph{lengste} adresseprefikset som matcher destinasjonsadressen.

\medskip
\textbf{Eksempel fra forelesning:}

\begin{center}
\begin{tabular}{lc}
\toprule
\textbf{Destinasjonsadresseomr\aa de} & \textbf{Link Interface} \\
\midrule
\texttt{11001000 00010111 00010***} & 0 \\
\texttt{11001000 00010111 00011000} & 1 \\
\texttt{11001000 00010111 00011***} & 2 \\
otherwise & 3 \\
\bottomrule
\end{tabular}
\end{center}

For adressen \texttt{11001000 00010111 00011000 10101010}:
\begin{itemize}[noitemsep]
    \item Matcher b\aa de interface 1 (full match p\aa\ 24 bit) og interface 2 (21 bit)
    \item Velger interface \textbf{1} fordi den har lengst prefiks-match
\end{itemize}

\medskip
\textbf{Fra eksamen V2023 (Q4):} Oppgaven ga en forwarding-tabell med 5 CIDR-oppf\o ringer og spurte hvilken utgangsport 5 ulike IP-adresser ville bli videresendt til. L\o sningen krever bin\ae r konvertering og longest prefix match.

\medskip
\textbf{Implementering:} Ofte implementert med \textbf{TCAM} (Ternary Content Addressable Memory) -- gir oppslag i \emph{en} klokkesyklus, uavhengig av tabellst\o rrelse.
\end{exambox}

\subsection{Svitsjefabrikk (Switching Fabric)}

\begin{keybox}[Tre typer svitsjefabrikk]
Svitsjefabrikken overforer pakker fra inngangsport til riktig utgangsport.

\textbf{Switching rate:} Hastigheten pakker kan overfores fra inngang til utgang.
\begin{itemize}[noitemsep]
    \item Ofte m\aa lt som multippel av inngangs-/utgangs-linjehastighet
    \item Med $N$ inngangsporter: \o nskelig med switching rate $= N \times R$ (linjehastighet)
\end{itemize}

\textbf{Tre hovedtyper:}
\begin{enumerate}[noitemsep]
    \item \textbf{Via minne:} F\o rste generasjon. Pakke kopieres til systemminne via CPU. Begrenset av minneb\aa ndbredde (2 busskryssinger per datagram).
    \item \textbf{Via buss:} Datagram fra inngangsportens minne til utgangsportens minne via en delt buss. \emph{Bus contention} -- hastighet begrenset av bussb\aa ndbredde. F.eks. Cisco 5600: 32 Gbps.
    \item \textbf{Via kryssbarnett (interconnection network):} Crossbar, Clos-nettverk. Kan overf\o re flere pakker parallelt. Fragmenterer datagrammer til celler, svitsjer cellene, reassemblerer. Skalerer med parallelle switching ``planes'' (f.eks. Cisco CRS: 100s Tbps).
\end{enumerate}
\end{keybox}

\subsection{K\o er og forsinkelser i ruteren}

\subsubsection{Input port queuing}

\begin{keybox}[HOL-blokkering]
Hvis svitsjefabrikken er langsommere enn summen av inngangsportene, oppst\aa r k\o er ved inngangsportene.
\begin{itemize}[noitemsep]
    \item K\o forsinkelse og pakketap p\aa\ grunn av buffer-overflow
    \item \textbf{Head-of-the-Line (HOL) blocking:} Et datagram fremst i k\o en blokkerer andre datagrammer i samme k\o\ fra \aa\ komme videre, selv om deres utgangsport er ledig
\end{itemize}
\end{keybox}

\subsubsection{Output port queuing}

\begin{keybox}[Utgangsportkoing -- viktig!]
Buffring kreves n\aa r datagrammer ankommer fra svitsjefabrikken raskere enn linjeoverforingshastigheten.
\begin{itemize}[noitemsep]
    \item K\o forsinkelse og pakketap ved buffer-overflow
    \item \textbf{Drop policy:} Hvilke pakker kastes n\aa r bufferen er full?
    \item \textbf{Scheduling discipline:} Velger blant k\o ede datagrammer for sending
\end{itemize}
\end{keybox}

\subsubsection{Bufferdimensjonering}

\begin{formulabox}[Tommelregel for bufferstorrelse]
\textbf{RFC 3439:}
\[
B = RTT \times C
\]
der $RTT \approx 250$ ms (typisk) og $C$ = linkkapasitet.

\medskip
F.eks. $C = 10$ Gbps $\Rightarrow$ buffer $= 2{,}5$ Gbit.

\medskip
\textbf{Nyere anbefaling:} Med $N$ samtidige flyter:
\[
B = \frac{RTT \times C}{\sqrt{N}}
\]

\textbf{OBS:} For mye buffring kan ogs\aa\ v\ae re problematisk -- gir lange RTT-er (``bufferbloat''), d\aa rlig ytelse for sanntidsapplikasjoner, og tr\aa\ TCP-respons.
\end{formulabox}

\subsection{Buffer Management}

\begin{keybox}[Buffer management -- drop og marking]
\textbf{Drop-politikk} (n\aa r bufferen er full):
\begin{itemize}[noitemsep]
    \item \textbf{Tail drop:} Kast ankommende pakke (enkleste)
    \item \textbf{Priority drop:} Kast basert p\aa\ prioritet
\end{itemize}

\textbf{Marking} (signal om metning):
\begin{itemize}[noitemsep]
    \item ECN (Explicit Congestion Notification)
    \item RED (Random Early Detection)
\end{itemize}
\end{keybox}

\subsection{Pakkeplanlegging (Scheduling)}

\begin{exambox}[Eksamensrelevant: Kodisipliner (V2023 Q1.3.5)]
Fire vanlige scheduling-metoder:

\begin{enumerate}[noitemsep]
    \item \textbf{FCFS / FIFO:} F\o rst inn, f\o rst ut. Enkleste metode.
    \item \textbf{Priority scheduling:} Trafikk klassifiseres, k\o es etter klasse. Send alltid fra h\o yeste prioritetsk\o\ f\o rst. FCFS innenfor hver klasse.
    \item \textbf{Round Robin (RR):} Trafikk klassifiseres. Serveren syklisk skanner klassek\o ene, sender \emph{en} pakke fra hver klasse om gangen.
    \item \textbf{Weighted Fair Queuing (WFQ):} Generalisert Round Robin. Hver klasse $i$ har vekt $w_i$, og f\aa r andel:
    \[
    \frac{w_i}{\sum_j w_j}
    \]
    Gir minimum b\aa ndbreddegaranti per trafikkklasse.
\end{enumerate}

\medskip
\textbf{Fra eksamen V2023 (Q1.3.5):} ``A FIFO queue holds $k$ packets when a new packet arrives. Considering queuing delay only, how many packets must be transmitted before the newly arrived packet can be transmitted?''

Svar: $k$ pakker (alternativ c). Den nye pakken m\aa\ vente til alle $k$ pakkene foran i k\o en er sendt.
\end{exambox}

%=============================================================================
\section{Internet Protocol (IP) (4.3)}
%=============================================================================

\subsection{IPv4-datagramformat}

\begin{keybox}[IPv4-datagrammets header]
IPv4-headeren er normalt \textbf{20 bytes} (uten opsjoner):

\begin{center}
\begin{tabular}{ll}
\toprule
\textbf{Felt} & \textbf{Beskrivelse} \\
\midrule
Version (4 bit) & IP-versjon (4 for IPv4) \\
Header length (4 bit) & Headerl\ae ngde i 32-bits ord \\
Type of service (8 bit) & Differensiert tjeneste (DSCP, ECN) \\
Total length (16 bit) & Total datagramlengde i bytes \\
Identifier, flags, offset & Brukes for fragmentering/reassemblering \\
TTL (8 bit) & Time to live -- maks antall hopp \\
Upper layer protocol (8 bit) & Protokoll i payload (6=TCP, 17=UDP) \\
Header checksum (16 bit) & Kontrollsum for headeren \\
Source IP (32 bit) & Avsenders IP-adresse \\
Destination IP (32 bit) & Mottakers IP-adresse \\
Options (variabel) & Tidsstempel, record route, etc. \\
\bottomrule
\end{tabular}
\end{center}
\end{keybox}

\begin{formulabox}[Overhead]
\textbf{Minimum overhead} for TCP over IP:
\[
\text{Overhead} = \underbrace{20 \text{ bytes}}_{\text{IP-header}} + \underbrace{20 \text{ bytes}}_{\text{TCP-header}} = 40 \text{ bytes}
\]
Dette betyr at sm\aa\ applikasjonsdata f\aa r stor relativ overhead.
\end{formulabox}

\subsection{IP-adressering}

\begin{keybox}[IP-adresser -- grunnleggende]
\begin{itemize}[noitemsep]
    \item En IP-adresse er \textbf{32 bit} = 4 bytes, skrevet i ``dotted decimal''-notasjon
    \item Eksempel: \texttt{223.1.1.1} = \texttt{11011111 00000001 00000001 00000001}
    \item IP-adresser tilh\o rer \textbf{grensesnitt} (interfaces), ikke verter
    \item En ruter har typisk flere grensesnitt, og dermed flere IP-adresser
    \item En vert har normalt ett eller to grensesnitt (kabel og/eller WiFi)
\end{itemize}
\end{keybox}

\subsection{Subnett}

\begin{exambox}[Eksamensrelevant: Subnett og CIDR (V2025, V2023)]
\textbf{Subnett:} En gruppe enheter som kan n\aa\ hverandre uten \aa\ g\aa\ gjennom en ruter.

\textbf{CIDR} (Classless InterDomain Routing):
\begin{itemize}[noitemsep]
    \item Adresseformat: \texttt{a.b.c.d/x}, der \texttt{x} er antall bit i nettverksdelen (subnettprefikset)
    \item De f\o rste \texttt{x} bitene = nettverksdel (subnet part)
    \item De resterende $32 - x$ bitene = vertsdel (host part)
\end{itemize}

\medskip
\textbf{Beregning av antall adresser og verter:}
\begin{itemize}[noitemsep]
    \item Totalt antall adresser i et subnett: $2^{32-x}$
    \item \textbf{Brukbare vertsadresser:} $2^{32-x} - 2$ (trekk fra nettverksadresse og kringkastingsadresse)
\end{itemize}

\medskip
\textbf{Kringkastingsadresse (broadcast):} Alle vertsbiter satt til 1.

\medskip
\textbf{Fra eksamen V2025:}
\begin{itemize}[noitemsep]
    \item ``En organisasjon f\aa r tildelt adresseblokken \texttt{208.27.1.0/22} og deler den i 12 subnett. Hvor mange brukbare vertsadresser?''
    \item L\o sning: /22 gir $2^{10} = 1024$ adresser. 12 subnett krever $\lceil\log_2 12\rceil = 4$ bit, s\aa\ hvert subnett f\aa r $2^{10-4} = 2^6 = 64$ adresser, dvs. $64 - 2 = \textbf{62}$ brukbare vertsadresser.
    \item ``Hva er kringkastingsadressen for \texttt{172.18.16.0/21}?''
    \item L\o sning: /21 betyr 21 nettverksbiter, 11 vertsbiter. Sett alle vertsbiter til 1: \texttt{172.18.23.255}
\end{itemize}
\end{exambox}

\begin{warnbox}[Slik beregner du kringkastingsadressen]
\textbf{Steg for steg:}
\begin{enumerate}[noitemsep]
    \item Skriv nettverksadressen bin\ae rt
    \item Identifiser vertsdelen (de siste $32 - x$ bitene)
    \item Sett \emph{alle} vertsbiter til 1
    \item Konverter tilbake til desimal
\end{enumerate}

\textbf{Eksempel:} \texttt{172.18.16.0/21}
\begin{itemize}[noitemsep]
    \item Bin\ae rt: \texttt{10101100.00010010.000\textbf{10000.00000000}}
    \item Vertsdel (11 bit): de siste 11 bitene
    \item Sett til 1: \texttt{10101100.00010010.000\textbf{10111.11111111}}
    \item Desimal: \texttt{172.18.23.255}
\end{itemize}
\end{warnbox}

\subsection{IP-adresser: bin\ae r konvertering}

\begin{exambox}[Eksamensrelevant: Binær konvertering (V2023 Q1.3.1)]
\textbf{Fra eksamen V2023:} ``What is \texttt{193.32.216.9} in binary?''

L\o sning: Konverter hvert oktett separat:
\begin{itemize}[noitemsep]
    \item $193 = 128 + 64 + 1 = $ \texttt{11000001}
    \item $32 = $ \texttt{00100000}
    \item $216 = 128 + 64 + 16 + 8 = $ \texttt{11011000}
    \item $9 = 8 + 1 = $ \texttt{00001001}
\end{itemize}
Svar: \texttt{11000001.00100000.11011000.00001001}
\end{exambox}

\subsection{Subnettildeling fra CIDR-blokk}

\begin{exambox}[Eksamensrelevant: Subnettildeling (V2023 Q1.3.2)]
\textbf{Fra eksamen V2023:} ``An organization is granted the address block \texttt{193.32.216.0/22}. Which of the following addresses belong to this block?''

\textbf{L\o sningsmetode:}
\begin{enumerate}[noitemsep]
    \item Bestem nettverksdelen: de f\o rste 22 bitene
    \item Alle adresser med samme 22-bit prefiks tilh\o rer blokken
    \item Adresseomr\aa det er \texttt{193.32.216.0} -- \texttt{193.32.219.255}
\end{enumerate}
\end{exambox}

\subsection{Hierarkisk adressering og ruteaggregering}

\begin{keybox}[Ruteaggregering]
\textbf{Hierarkisk adressering} muliggj\o r effektiv annonsering av ruteinformasjon:
\begin{itemize}[noitemsep]
    \item En ISP f\aa r tildelt en stor adresseblokk (f.eks. \texttt{200.23.16.0/20})
    \item ISP-en deler blokken i mindre deler til organisasjoner (f.eks. 8 blokker \`a\ /23)
    \item Overfor resten av Internett annonserer ISP-en kun \emph{en} samlet rute (\texttt{200.23.16.0/20})
    \item Dette reduserer st\o rrelsen p\aa\ rutingtabeller dramatisk
\end{itemize}

\textbf{Mer spesifikke ruter:} Hvis en organisasjon bytter ISP, kan den nye ISP-en annonsere en mer spesifikk rute (f.eks. /23). P\aa\ grunn av longest prefix match vil trafikken g\aa\ til riktig ISP.
\end{keybox}

\subsection{Hvordan f\aa r en vert sin IP-adresse?}

\begin{keybox}[To metoder]
\begin{itemize}[noitemsep]
    \item \textbf{Manuell konfigurasjon:} Systemadministrator konfigurerer adressen direkte
    \item \textbf{DHCP (Dynamic Host Configuration Protocol):} Automatisk tildeling fra server
\end{itemize}
\end{keybox}

\subsection{DHCP -- Dynamic Host Configuration Protocol}

\begin{exambox}[Eksamensrelevant: DHCP (V2025, V2023)]
DHCP er en ``plug-and-play''-protokoll som lar en vert automatisk f\aa\ en IP-adresse.

\textbf{DHCP bruker UDP:} Klient p\aa\ port 68, server p\aa\ port 67.

\textbf{Fire-stegs prosess (DORA):}
\begin{enumerate}[noitemsep]
    \item \textbf{DHCP Discover:} Klienten sender broadcast (\texttt{255.255.255.255}) for \aa\ finne DHCP-servere.
    \begin{itemize}[noitemsep]
        \item SIP=\texttt{0.0.0.0}, DIP=\texttt{255.255.255.255}
        \item Ethernet DA=\texttt{FF:FF:FF:FF:FF:FF}
    \end{itemize}

    \item \textbf{DHCP Offer:} Server(e) svarer med tilbudt IP-adresse, subnettmaske, DNS, ruter, lease-tid.

    \item \textbf{DHCP Request:} Klienten velger et tilbud og broadcaster en foresp\o rsel med server-ID.
    \begin{itemize}[noitemsep]
        \item SIP=\texttt{0.0.0.0} (fortsatt), DIP=\texttt{255.255.255.255} (broadcast)
        \item Inneholder \texttt{DHCP Server Identifier} for valgt server
        \item Broadcast sikrer at andre servere vet at de \emph{ikke} ble valgt
    \end{itemize}

    \item \textbf{DHCP ACK:} Valgt server bekrefter tildelingen med alle konfigurasjonsparametere.
\end{enumerate}

\medskip
\textbf{DHCP gir mer enn bare IP-adresse:}
\begin{itemize}[noitemsep]
    \item Adresse til f\o rstehopp-ruter (default gateway)
    \item Navn og IP-adresse til DNS-server
    \item Nettverksmaske (subnettmaske)
\end{itemize}

\medskip
\textbf{Protokollstabel:} DHCP-melding kapsles inn i UDP $\rightarrow$ IP $\rightarrow$ Ethernet.

\medskip
\textbf{Fra eksamen V2023 (Q1.3.3):} ``Which of these are related to DHCP?''
\begin{itemize}[noitemsep]
    \item a) Provides an IP address to a host -- \textbf{korrekt}
    \item b) Uses UDP -- \textbf{korrekt}
    \item d) Address of default gateway is provided -- \textbf{korrekt}
\end{itemize}
\end{exambox}

\subsection{NAT -- Network Address Translation}

\begin{keybox}[NAT -- oversikt]
NAT lar flere enheter i et lokalt nettverk dele \emph{en} offentlig IP-adresse.

\textbf{Hvordan det fungerer:}
\begin{itemize}[noitemsep]
    \item Alle datagrammer som forlater lokalnettet f\aa r \emph{samme} kilde-IP (NAT-ruterens offentlige adresse), men \emph{ulike} kildeportnumre
    \item NAT-ruteren har en \textbf{oversettingstabell} (NAT translation table) som mapper: \\
    (kilde-IP, kildeport) $\leftrightarrow$ (NAT-IP, NAT-port)
    \item Innkommende pakker oversettes tilbake ved hjelp av tabellen
\end{itemize}

\textbf{Private adresseomr\aa der} (RFC 1918):
\begin{itemize}[noitemsep]
    \item \texttt{10.0.0.0/8} -- 16.7 mill. adresser
    \item \texttt{172.16.0.0/12} -- ca. 1 mill. adresser
    \item \texttt{192.168.0.0/16} -- 65\,536 adresser
\end{itemize}

\textbf{16-bit portnumre} gir mulighet for ca. 65\,000 samtidige forbindelser med en enkelt IP-adresse.
\end{keybox}

\begin{exambox}[Eksamensrelevant: NAT og private adresser (V2025)]
\textbf{Fra eksamen V2025 (sant/usant):}
\begin{itemize}[noitemsep]
    \item ``IP-adresser er alltid globalt unike'' -- \textbf{USANT}. Private adresser (10.x.x.x, 172.16.x.x, 192.168.x.x) gjenbrukes i mange nettverk, og er kun unike innenfor sitt lokalnett. NAT er grunnen til at dette fungerer.
\end{itemize}
\end{exambox}

\begin{warnbox}[NAT-kontrovers]
NAT er kontroversielt fordi:
\begin{itemize}[noitemsep]
    \item Rutere skal kun operere opp til lag 3 -- NAT manipulerer portnumre (lag 4)
    \item Bryter ende-til-ende-prinsippet
    \item Gj\o r det vanskelig for eksterne verter \aa\ initialisere kontakt med enheter bak NAT
    \item IPv6 var ment \aa\ gj\o re NAT un\o dvendig (nok adresser til alle)
\end{itemize}
\end{warnbox}

\subsection{IPv6}

\begin{keybox}[IPv6 -- hovedforskjeller fra IPv4]
\textbf{Motivasjon:} IPv4-adresserommet (32 bit = ca. 4.3 mrd.) er oppbrukt. ICANN tildelte siste blokk i 2011.

\textbf{IPv6-datagramformat:}
\begin{itemize}[noitemsep]
    \item \textbf{128-bit} adresser (vs. 32-bit i IPv4)
    \item Fast headerstorrelse p\aa\ \textbf{40 bytes} (vs. variabel i IPv4)
    \item \textbf{Ingen fragmentering} av rutere -- kun endesystemer kan fragmentere
    \item \textbf{Ingen header checksum} -- raskere prosessering
    \item \textbf{Nytt felt: Flow Label} -- identifiserer pakker i samme flyt for spesialbehandling
    \item \texttt{Next Header}-felt erstatter b\aa de ``Protocol'' og ``Options'' fra IPv4
    \item ICMPv6: ny versjon av ICMP med tilleggsmeldingstyper
\end{itemize}
\end{keybox}

\begin{exambox}[Eksamensrelevant: IPv4 vs. IPv6 (V2023 Q1.3.4)]
\textbf{Fra eksamen V2023:} ``Which of these are correct about IPv4/IPv6?''
\begin{itemize}[noitemsep]
    \item a) IPv6 has a larger address space -- \textbf{korrekt} (128 vs. 32 bit)
    \item b) IPv6 has a fixed-length header -- \textbf{korrekt} (40 bytes)
    \item d) IPv6 does not allow fragmentation by routers -- \textbf{korrekt}
    \item e) IPv6 removes the header checksum field -- \textbf{korrekt}
\end{itemize}

\textbf{Fra eksamen V2025 (sant/usant):}
\begin{itemize}[noitemsep]
    \item ``IP guarantees that a packet will be delivered in order'' -- \textbf{USANT}. IP er ``best effort'' og gir ingen garanti for rekkef\o lge. Det er TCPs jobb.
\end{itemize}
\end{exambox}

\begin{keybox}[Felt som er fjernet i IPv6 (sammenlignet med IPv4)]
\begin{itemize}[noitemsep]
    \item \textbf{Header checksum:} Fjernet for raskere prosessering (kontrollsum beregnes allerede i transport- og linklaget)
    \item \textbf{Fragmentering/reassemblering:} Rutere fragmenterer ikke -- avsender m\aa\ bruke Path MTU Discovery
    \item \textbf{Options:} Erstattet av ``Next Header''-mekanismen
    \item \textbf{Header Length:} Un\o dvendig med fast headerstorrelse
\end{itemize}
\end{keybox}

\subsection{Overgang fra IPv4 til IPv6: Tunneling}

\begin{keybox}[Tunneling -- IPv6 i IPv4]
Siden ikke alle rutere kan oppgraderes samtidig, brukes \textbf{tunneling}:
\begin{itemize}[noitemsep]
    \item IPv6-datagrammet pakkes inn som \emph{payload} i et IPv4-datagram
    \item IPv4-rutere langs tunnelen ser det som en vanlig IPv4-pakke
    \item Ved tunnelens ende pakkes IPv6-datagrammet ut igjen
\end{itemize}

\textbf{Adopsjon:} Ca. 40--50\% av trafikken bruker IPv6 (per ca. 2023).
\end{keybox}

%=============================================================================
\section{Arkitekturprinsipper (4.5)}
%=============================================================================

\subsection{IP-timeglass-arkitekturen}

\begin{keybox}[IP som ``smal midje'']
Internett-arkitekturen har en timeglass-form:
\begin{itemize}[noitemsep]
    \item \textbf{Over IP:} Mange applikasjoner og transportprotokoller (HTTP, SMTP, TCP, UDP, \ldots)
    \item \textbf{IP:} Det eneste nettverkslagset -- ``smal midje'' som alt m\aa\ g\aa\ gjennom
    \item \textbf{Under IP:} Mange linklag-teknologier (Ethernet, WiFi, fiber, \ldots)
\end{itemize}
Alle enheter p\aa\ Internett m\aa\ st\o tte IP -- det er det som binder alt sammen.
\end{keybox}

\subsection{Middleboxes}

\begin{keybox}[Middleboxes]
Enheter p\aa\ nettverksveien som utf\o rer funksjoner utover standard IP-forwarding:
\begin{itemize}[noitemsep]
    \item \textbf{NAT-bokser:} Adresseoversetting
    \item \textbf{Brannmurer:} Filtrering av trafikk
    \item \textbf{Lastbalanserere:} Fordeler trafikk over flere servere
    \item \textbf{Cacher/proxyer:} Mellomlagring av innhold
\end{itemize}

Middleboxes er omstridte fordi de potensielt bryter med \textbf{ende-til-ende-prinsippet}.
\end{keybox}

\subsection{Ende-til-ende-argumentet}

\begin{keybox}[End-to-End Argument]
\textbf{Ende-til-ende-prinsippet:} Funksjonalitet b\o r implementeres i \emph{endesystemene} (vertene), ikke i nettverket.

\begin{itemize}[noitemsep]
    \item Nettverket b\o r holdes \textbf{enkelt} -- kun videresending av pakker
    \item Kompleksitet (p\aa litelighet, sikkerhet, osv.) b\o r ligge i endesystemene
    \item Begrunnelse: Endesystemene har mest kontekst om hva applikasjonen trenger
\end{itemize}

\textbf{Historisk utvikling:}
\begin{itemize}[noitemsep]
    \item 20. \aa rhundre (telefonnett): Intelligens i nettverkets svitsjer
    \item Internett (pre-2005): Intelligens i ytterkantene (endesystemer)
    \item Internett (post-2005): Programmerbare nettverksenheter + massiv infrastruktur i ytterkantene
\end{itemize}
\end{keybox}

%=============================================================================
\section{IP-adressering: ICANN og adressetildeling}
%=============================================================================

\begin{keybox}[Hvordan f\aa r en ISP adresser?]
\textbf{ICANN} (Internet Corporation for Assigned Names and Numbers):
\begin{itemize}[noitemsep]
    \item Tildeler IP-adresseblokker gjennom 5 \textbf{regionale registre} (RIR-er)
    \item Administrerer DNS-rotsonen
    \item Delegerer TLD-administrasjon (.com, .no, .edu, osv.)
\end{itemize}

\textbf{Adressetildelingshierarki:}
\begin{enumerate}[noitemsep]
    \item ICANN $\rightarrow$ Regionale registre (RIPE NCC for Europa)
    \item Regionale registre $\rightarrow$ ISP-er
    \item ISP-er $\rightarrow$ Organisasjoner/kunder (via CIDR-blokker)
    \item Organisasjoner $\rightarrow$ Enkelthoster (via DHCP eller manuelt)
\end{enumerate}
\end{keybox}

%=============================================================================
\section{Oppsummering og eksamenstips}
%=============================================================================

\begin{exambox}[Ofte stilte eksamenstemaer -- Kapittel 4]
Basert p\aa\ eksamener fra V2023 og V2025:

\begin{enumerate}[noitemsep]
    \item \textbf{Forwarding vs. Routing:} Klar definisjon og skille mellom lokal/global, dataplanet/kontrollplanet
    \item \textbf{Subnettberegning:} Kunne beregne antall verter, kringkastingsadresse, sjekke om en adresse tilh\o rer et subnett (CIDR)
    \item \textbf{Bin\ae r konvertering:} Konvertere mellom desimal og bin\ae r IP-notasjon
    \item \textbf{Longest prefix match:} Sl\aa\ opp i forwarding-tabell med CIDR-oppf\o ringer
    \item \textbf{DHCP:} 4-stegs prosessen (DORA), bruk av broadcast, UDP-porter, hva DHCP gir
    \item \textbf{IPv4 vs. IPv6:} Forskjeller i header, adresselengde, fragmentering, checksum
    \item \textbf{NAT:} Hvordan det fungerer, private adresser, konsekvenser for ``globalt unike'' adresser
    \item \textbf{K\o disipliner:} FIFO, priority, round robin, WFQ -- spesielt FIFO-beregninger
    \item \textbf{Best effort:} IP gir ingen garantier (rekkef\o lge, levering, forsinkelse)
\end{enumerate}
\end{exambox}

\begin{warnbox}[Vanlige feller p\aa\ eksamen]
\begin{itemize}[noitemsep]
    \item \textbf{``IP garanterer rekkef\o lge''} -- \textbf{Feil!} IP er best effort, ingen rekkef\o lgegaranti
    \item \textbf{``IP-adresser er alltid globalt unike''} -- \textbf{Feil!} Private adresser (NAT) gjenbrukes
    \item \textbf{Forwarding $\neq$ Routing} -- Ikke bland disse begrepene
    \item \textbf{Glem ikke \aa\ trekke fra 2} n\aa r du beregner brukbare vertsadresser (nettverksadresse + broadcast)
    \item \textbf{Broadcast-adresse:} Husk \aa\ sette \emph{alle} vertsbiter til 1, ikke bare siste oktett
    \item \textbf{DHCP Request er broadcast}, ikke unicast (s\aa\ andre servere vet at de ikke ble valgt)
\end{itemize}
\end{warnbox}

\begin{formulabox}[Viktige formler -- Kapittel 4]
\begin{align}
\text{Antall adresser i subnett } /x &= 2^{32-x} \\
\text{Brukbare vertsadresser} &= 2^{32-x} - 2 \\
\text{Antall subnett-biter for } n \text{ subnett} &= \lceil \log_2 n \rceil \\
\text{Bufferstorrelse (RFC 3439)} &= RTT \times C \\
\text{Bufferstorrelse (med } N \text{ flyter)} &= \frac{RTT \times C}{\sqrt{N}} \\
\text{WFQ-andel for klasse } i &= \frac{w_i}{\sum_j w_j} \\
\text{IP+TCP overhead} &= 20 + 20 = 40 \text{ bytes}
\end{align}
\end{formulabox}
